\begin{resumen}%
En este trabajo implementamos dos sistemas para medir los pulsos de un detector
sobre la Red Pitaya STEMlab 125-14, una placa con un SoC Zynq-7010 y un ADC de
14~bits a 125~MS/s. El osciloscopio por eventos (OSC) captura la traza completa
de cada pulso, para analizarla después. El analizador multicanal (MCA), la
herramienta central de la espectroscopía nuclear, no guarda los pulsos sino sus
estadísticos. Los dos corren en la
lógica programable (PL) y comparten el disparo, con seis fuentes combinables.
El MCA mide la amplitud con un banco de tres filtros trapezoidales y la forma
del pulso, que dan un espectro de 8192 canales y un mapa de amplitud $\times$
forma, y usa tres ranuras de medición en paralelo para reducir el tiempo
muerto. El OSC escribe las trazas en un anillo de eventos en la DDR. El
procesador (PS) se queda con el control, la
adquisición y el guardado, a través de una API en C y Python, un servidor de red
y una aplicación gráfica. El diseño ocupa toda la FPGA y cierra el \emph{timing} a
\qty{125}{\mega\hertz}.

Para caracterizar el sistema agregamos en la PL un generador de pulsos con
alturas, formas, ruido y arribos programables, y medimos en el lazo digital y en
el canal analógico de la placa. En el lazo digital la cadena no agrega un error
medible: la INL del trapecio queda por debajo de \qty{0.002}{\percent}~FS y las
líneas se corren menos de 1.5 canales hasta 3~M/s. Con 3~M/s de entrada el MCA
acepta 1.9~M/s, y con el disparo por polo-cero la tasa estimada queda a
\qty{3.7}{\percent} de la real; lo que limita la tasa ya no es el tiempo muerto
sino el apilamiento de los pulsos. En el canal analógico lo que limita es su
respuesta al escalón: el sobrepico y la caída lenta. El OSC guarda desde Python
100~k/s trazas de 32 muestras, limitado por el PS y no por la PL ni por la
memoria. El MCA no tiene esa limitación, porque acumula el histograma en la PL y
el PS sólo lee el espectro.

Por último, proponemos los pasos que quedan para llevar el sistema a una
medición con un detector real: compensar la caída lenta del canal, estabilizar
el espectro con la temperatura y compararlo con un MCA comercial.
\end{resumen}



%%% Local Variables:
%%% mode: latex
%%% TeX-master: "template"
%%% End:
